\scanpage{113}
\mspage{62}
sous-topos fermé complémentaire, est toute similaire.
(\uline{NB} Si $\unsure{\mathscr X|\mathcal X}=\mathrm{Top}(X)$, avec $X$ dans $\mathrm{Cat}$, alors $\mathrm{Ouv}(\mathscr X)\simeq\deleted{\unsure{\mathrm{Ouv}}}\mathrm{Ouv}(X)$, puisque les ouverts correspondent aux sous-faisceaux du faisceau final sur $\mathscr X$, i.e. du préfaisceau final sur $X$. Cette bijection est compatible avec le passage aux complémentaires.)
La théorie des topos dit que les couples $(\mathscr X,\mathscr U,\unsure{\mathscr T|\mathcal T})$ se réconstituent, à \unsure{équiv.|équi.} définie à isom unique près, à l'aide d'un triple
\[
\left\{\begin{aligned}&\mathscr U,\mathscr T&&\text{deux topos}\\ &\varphi:\mathrm{Fais}(\mathscr U)\longrightarrow\mathscr T &&\text{un foncteur exact}\\ &&&\text{à gauche (``accessible''),}\end{aligned}\right.
\]
ce foncteur $\varphi$ étant obtenu, quand on part de
\diagram{\begin{tikzcd}[ampersand replacement=\&]
\mathscr U\arrow[r,"i"] \& \mathscr X \& \mathscr T\arrow[l,"j"'] ,
\end{tikzcd}}
comme
\[
\varphi=j^*i_*.
\]
De plus, ce triple $(\mathscr U,\mathscr T,\varphi)$ peut être choisi arbitrairement, avec la seule restriction (pratiquement anodine, car toujours satisfaite pour tous les foncteurs naturels qu'on rencontre) que $\varphi$ soit (en plus d'être exact à gauche) ``accessible''.

Dans le cas actuel \deleted{\unsure{$\mathscr U$|$\mathscr X$}} $\mathscr X=\mathrm{Top}\,X,\mathscr U=\mathrm{Top}\,U,\mathscr T=\mathrm{Top}(T)$, le foncteur \deleted{\unsure{$\mathrm{Top}$}} $\varphi=j^*i_*$ \deleted{$\mathscr X$ est} non seulement exact à gauche, mais même il commute aux $\varprojlim$ quelconques, puisque il en est ainsi (toujours) de $i_*$, et aussi (puisque \inserted{le} $j$ \inserted{\unsure{toposique}} est induit par le $j:T\longrightarrow X$ catégorique) de $j^*$. En fait on
% archive sheet number: 113
